\chapter{符号、约定与证明依赖}
除特别说明，流形指有限维、Hausdorff、第二可数的实光滑流形，无边界；带边情形单独声明。张量型别采用$(r,s)$表示$r$个逆变指标与$s$个协变指标。重复的上下指标求和。
\section{曲率与符号}
统一采用
\[
R(X,Y)Z=\nabla_X\nabla_YZ-\nabla_Y\nabla_XZ-\nabla_{[X,Y]}Z,
\qquad
K(X,Y)=\frac{g(R(X,Y)Y,X)}{g(X,X)g(Y,Y)-g(X,Y)^2}.
\]
这使标准单位球面的截面曲率为正。定义$\Ric(Y,Z)=\operatorname{tr}(X\mapsto R(X,Y)Z)$。Lorentz度量采用$(-,+,+,+)$；物理常数、单位及Einstein方程的系数须在相对论正文中明确。不同参考书的符号不能直接混用。
\section{分析工具的证明边界}
逆映射、隐映射、常秩、常微分方程局部存在唯一性及光滑依赖的完整陈述和证明列入预备知识待办。Hodge理论与指标理论还需要Sobolev空间、椭圆正则性和Fredholm理论。框架版不宣称这些证明已经具备，也不在缺少这些工具时宣称推得相应结论。
